
\documentclass{abgabe}
\begin{document}

\def\🚗{\scalerel*{\includegraphics{figures/1f697.eps}}{\textrm{\textbigcircle}}}

\begin{questions}
    \question
    Ein Bahnübergang wird in der Zeit zwischen 16.00 Uhr und 18.00 Uhr durchschnittlich von 150 Fahrzeugen pro Stunde überquert. 
    Bestimmen Sie die Wahrscheinlichkeit, dass während der jeweils für $90\si{\s}$ geschlossenen Schranken
    \begin{parts}
        \part 
        kein Fahrzeug 
        \begin{solution}
            Es gilt offensichtlich: 
            \[ 
                \frac{150 \ \🚗}{1 \si{\hour}} = \frac{150 \ \🚗}{3600 \si{\s}} = \frac{\nicefrac{15}{4} \ \🚗}{90 \si{\s}} 
            \] 
            Damit erhalten wir unseren Erwartungswert $\mu \ [ \🚗 ]$ mit $\mu = \frac{15}{4}$.
            Offenbar gilt: 
            \[ 
                X = \{ \text{Anzahl der Autos in 90\si{\s}} \} \sim \po(x;\mu) = \po(x;\nicefrac{15}{4})
            \] 
            mit 
            \[ 
                P(X = x) = \frac{\mu^x}{x!} \cdot e^{-\mu} = \frac{\left( \nicefrac{15}{4} \right)^x}{x!} \cdot e^{-\nicefrac{15}{4}}
            \] 
            
            Und damit:
            \[ 
            P(X = 0) = \frac{\left( \nicefrac{15}{4} \right)^0}{0!} \cdot e^{-\nicefrac{15}{4}} \approx 2,35\%
            \] 
        \end{solution}
        
        \part 
        genau ein Fahrzeug 
        \begin{solution}
            Es gilt:
            \[ 
            P(X = 1) = \frac{\left( \nicefrac{15}{4} \right)^1}{1!} \cdot e^{-\nicefrac{15}{4}} \approx 8.82\%
            \] 
        \end{solution}
        
        \part 
        genau zwei Fahrzeuge
        \begin{solution}
            Es gilt:
            \[ 
            P(X = 2) = \frac{\left( \nicefrac{15}{4} \right)^2}{2!} \cdot e^{-\nicefrac{15}{4}} \approx 16.54\%
            \] 
        \end{solution}
        
        \part 
        mehr als zwei Fahrzeuge
        \begin{solution}
            Es gilt: 
            \[ 
                P(X > 2) = 1 - P(X \leq 2) = 1 - (P(X=0) + P(X=2) + P(X=2)) \approx 72.29\% 
            \] 
        \end{solution}
    \end{parts}
    eintreffen.
\end{questions}
\end{document}